\documentclass[conference]{IEEEtran}
\IEEEoverridecommandlockouts

\usepackage{cite}
\usepackage{amsmath,amssymb,amsfonts}
\usepackage{graphicx}
\usepackage{textcomp}
\usepackage{xcolor}
\usepackage{booktabs}
\usepackage{listings}
\usepackage{url}
\usepackage{hyperref}

\lstset{
  basicstyle=\ttfamily\footnotesize,
  breaklines=true,
  frame=single
}

\begin{document}

\title{SIMD Vector Processing Unit: Design, Verification, and FPGA Synthesis}

\author{
  \IEEEauthorblockN{Kasula Lalithendra}
  \IEEEauthorblockA{\textit{B.Tech CS\&AI, Rishihood University / NST}\\
  kasula.lalithendra2024@nst.rishihood.edu.in}
  \and
  \IEEEauthorblockN{Dhanvin Vadlamudi}
  \IEEEauthorblockA{\textit{B.Tech CS\&AI, Rishihood University / NST}\\
  dhanvin.vadlamudi269@gmail.com}
  \and
  \IEEEauthorblockN{Jagruthi Pulumati}
  \IEEEauthorblockA{\textit{B.Tech CS\&AI, Rishihood University / NST}\\
  jagruthip2007@gmail.com}
  \and
  \IEEEauthorblockN{Rakshita}
  \IEEEauthorblockA{\textit{B.Tech CS\&AI, Rishihood University / NST}\\
  polanarakshita28@gmail.com}
  \and
  \IEEEauthorblockN{Chaitanya Sai Meka}
  \IEEEauthorblockA{\textit{B.Tech CS\&AI, Rishihood University / NST}\\
  chaitanya.s.meka@gmail.com}
}

\maketitle

% ─────────────────────────────────────────────────────────────────────────────
\begin{abstract}
This paper presents the design, implementation, verification, and FPGA synthesis
of a parameterized SIMD (Single Instruction, Multiple Data) Vector Processing
Unit (VPU) in Verilog-2001. The VPU executes one instruction simultaneously
across \texttt{LANES} parallel 32-bit data lanes (default: 4~lanes), supporting
nine ALU operations (ADD, SUB, AND, OR, XOR, SLT, EQ, SLL, SRL), vector
load/store, scalar-broadcast addressing, and a minimal scalar control interface.
A non-pipelined Harvard-architecture FSM controls instruction fetch, decode/execute,
and write-back in deterministic cycle counts. We verify the RTL against a Python
golden model using self-checking testbenches for every module, and synthesize
the design for Artix-7 (\texttt{xc7a35tcpg236-1}) at lane counts 1, 2, 4, and~8.
Measured results include cycle count, throughput, LUT/FF/DSP resource usage,
and maximum operating frequency. The 4-lane configuration achieves approximately
$4\times$ throughput improvement over the scalar (LANES\,=\,1) baseline on pure
data-parallel benchmarks, with sub-linear scaling losses attributable to loop
overhead and scalar control instructions.
% TODO: fill in actual Fmax and resource numbers from synth/reports/ before submission
\end{abstract}

\begin{IEEEkeywords}
SIMD, vector processor, FPGA, Verilog, data-level parallelism, hardware design
\end{IEEEkeywords}

% ─────────────────────────────────────────────────────────────────────────────
\section{Introduction and Objectives}

Modern workloads in signal processing, machine learning, and scientific computing
exhibit abundant \emph{data-level parallelism} (DLP): the same operation is applied
independently to large arrays of elements. SIMD architectures exploit DLP by
issuing a single instruction that operates on multiple data elements in parallel,
amortising the fetch-decode overhead across all lanes~\cite{flynn1972some}.

This project designs and evaluates a small, synthesizable SIMD VPU as a learning
vehicle for key concepts in computer architecture: register files, ALU design,
parameterized RTL, Harvard memory, FSM-based control, and FPGA synthesis. The
primary objectives are:

\begin{enumerate}
  \item Implement a parameterized VPU in synthesizable Verilog-2001 that supports
        LANES\,$\in\{1, 2, 4, 8\}$.
  \item Verify all modules against a Python golden reference model using
        self-checking testbenches.
  \item Measure the cycle-count speedup of LANES\,=\,4 and LANES\,=\,8 over the
        scalar baseline (LANES\,=\,1) on three benchmarks.
  \item Synthesize for Artix-7 and report resource utilization and maximum
        frequency at all four lane counts.
\end{enumerate}

% ─────────────────────────────────────────────────────────────────────────────
\section{Literature Review}

\subsection{Flynn's Taxonomy and SIMD}

Flynn classified computer architectures by the multiplicity of instruction and
data streams~\cite{flynn1966very}. The SIMD class — Single Instruction, Multiple
Data — is the foundation of modern vector and SIMD hardware, from early array
processors to contemporary CPU extensions (SSE, AVX, NEON) and GPU shader cores.
Flynn later formalized the relationship between instruction throughput and
parallelism in~\cite{flynn1972some}.

\subsection{Data-Level Parallelism and Amdahl's Law}

Amdahl~\cite{amdahl1967validity} showed that the speedup of a parallel machine
is bounded by the fraction of the workload that is parallelizable. For a pure
vector workload (Amdahl fraction $f \to 1$) with $L$ lanes, the ideal speedup
approaches $L$. In practice, loop overhead (scalar address arithmetic, branch,
HALT) and memory bandwidth limit the achieved speedup. We measure this effect
directly in Section~VI.

\subsection{Vector Processor Architecture}

Hennessy and Patterson~\cite{hennessy2017computer} provide the standard reference
for vector processor design: vector register files, lane-parallel functional
units, vector load/store units, and stride-based addressing. Our design follows
the same principles at a small scale suitable for an FPGA prototype.

\subsection{FPGA Soft SIMD}

Reconfigurable devices allow SIMD datapaths to be implemented as soft processors.
% TODO: add one or two survey papers on FPGA soft processors or vector accelerators
% if available; mark with \% VERIFY in refs.bib

% ─────────────────────────────────────────────────────────────────────────────
\section{Proposed Solution / Architecture}

\subsection{Top-Level Architecture}

The VPU follows a Harvard architecture (separate instruction and data memories)
with a non-pipelined, single-issue FSM. Figure~\ref{fig:block} shows the
Phase~I demo datapath.

\begin{figure}[htbp]
  \centering
  % TODO: replace with actual exported PNG from docs/block_diagram.md
  \fbox{\rule{0pt}{4cm}\rule{0.9\columnwidth}{0pt}}
  \caption{SIMD VPU demo datapath block diagram (LANES\,=\,4). Signals flow from
  \texttt{vreg\_file} and \texttt{sreg\_file} through \texttt{lane\_array}
  (which replicates \texttt{lane\_mux} and \texttt{lane\_alu} for each lane)
  to \texttt{result\_combiner} and back to \texttt{vreg\_file}.}
  \label{fig:block}
\end{figure}

\subsection{Module Hierarchy}

Table~\ref{tab:modules} summarises the RTL modules. Full port lists and bus
conventions are in \texttt{docs/interfaces.md}.

\begin{table}[htbp]
\caption{RTL Module Summary}
\label{tab:modules}
\centering
\begin{tabular}{lp{5cm}}
\toprule
Module & Responsibility \\
\midrule
\texttt{lane\_alu}      & 32-bit single-lane ALU (9 ops, combinational) \\
\texttt{lane\_mux}      & Operand-B select: vector lane or scalar broadcast \\
\texttt{lane\_array}    & \texttt{LANES} copies via \texttt{generate} loop \\
\texttt{vreg\_file}     & 8 vector registers, async read, sync write \\
\texttt{sreg\_file}     & 8 scalar registers, \texttt{r0}$\equiv$0 \\
\texttt{result\_combiner} & ALU/load write-back mux \\
\texttt{dmem}           & Wide-port data memory (\texttt{LANES}$\times$32\,b) \\
\texttt{pc\_imem}       & Program counter and instruction memory \\
\texttt{vpu\_demo\_top} & Phase~I demo top (Phase~II adds \texttt{control\_unit}) \\
\bottomrule
\end{tabular}
\end{table}

\subsection{Instruction Set}

The VPU uses 32-bit fixed-width instructions. The instruction encoding
(opcode[31:26], destination[25:23], src1[22:20], src2[19:17],
opb\_sel[16], imm[15:0]) is documented in \texttt{docs/isa.md}.
Supported operations: VADD, VSUB, VAND, VOR, VXOR, VSLT, VEQ, VSLL,
VSRL, VLD, VST, SADDI, SLI, BNE, HALT.

\subsection{Control FSM}

The control unit (Phase~II) is a 3-state Moore FSM:
FETCH $\to$ DECODE/EXECUTE $\to$ WRITEBACK.
Vector ALU and VLD require 3~cycles; VST, BNE, and HALT require 2~cycles.
The FSM design is specified in \texttt{docs/control\_fsm.md}.

% ─────────────────────────────────────────────────────────────────────────────
\section{Individual Contributions}

\begin{table}[htbp]
\caption{Team Contributions (fill from \texttt{git log --author} before submission)}
\label{tab:contributions}
\centering
\begin{tabular}{lp{5.5cm}}
\toprule
Team Member & Contribution \\
\midrule
Kasula Lalithendra  & \\
Dhanvin Vadlamudi   & \\
Jagruthi Pulumati   & \\
Rakshita            & \\
Chaitanya Sai Meka  & \\
\bottomrule
\end{tabular}
\end{table}

% ─────────────────────────────────────────────────────────────────────────────
\section{Verification}

All modules are verified with self-checking testbenches compiled and run under
Icarus Verilog. The Python golden model (\texttt{verif/golden.py}) computes
expected outputs for all nine ALU operations across edge-case inputs
(0, 1, $-1$, \texttt{0x7FFFFFFF}, \texttt{0x80000000}, overflow wraparound,
shifts by 0 and~31). Testbenches print \texttt{PASS} with a count or
\texttt{FAIL: expected=X actual=Y} per check.

\texttt{lane\_array\_tb} is compiled three times at LANES\,$\in\{1,4,8\}$
and includes a lane-independence test that mutates one lane's input and
confirms only that lane's output changes. All tests pass prior to submission.
% TODO: paste final 'make test' output here before submission.

A waveform screenshot of a 4-lane VADD showing all four lanes updating in the
same clock cycle is in \texttt{docs/}.

% ─────────────────────────────────────────────────────────────────────────────
\section{Results}

% TODO: fill in after Vivado synthesis and benchmark runs complete
% Numbers must come from report/results.csv (make bench)

\subsection{Cycle Count and Speedup}

Table~\ref{tab:cycles} shows \texttt{cycle\_count} for benchmark~B1
(\texttt{C[i] = A[i] + B[i]}, $N = 64$) at each lane count.

\begin{table}[htbp]
\caption{B1 Benchmark Cycle Count and Speedup ($N=64$)}
\label{tab:cycles}
\centering
\begin{tabular}{rrrr}
\toprule
LANES & Cycles & Cycles/Element & Speedup vs.\ LANES\,=\,1 \\
\midrule
1 & \textit{TBD} & \textit{TBD} & 1.00$\times$ \\
2 & \textit{TBD} & \textit{TBD} & \textit{TBD} \\
4 & \textit{TBD} & \textit{TBD} & \textit{TBD} \\
8 & \textit{TBD} & \textit{TBD} & \textit{TBD} \\
\bottomrule
\end{tabular}
\end{table}

\subsection{FPGA Resource Utilization and Fmax}

Table~\ref{tab:synth} summarises Vivado synthesis results for the
\texttt{lane\_alu} module on \texttt{xc7a35tcpg236-1}.

\begin{table}[htbp]
\caption{Synthesis Results — \texttt{lane\_alu} on xc7a35tcpg236-1}
\label{tab:synth}
\centering
\begin{tabular}{rrrrrr}
\toprule
LANES & LUTs & FFs & DSPs & BRAMs & Fmax (MHz) \\
\midrule
1 & \textit{TBD} & \textit{TBD} & \textit{TBD} & \textit{TBD} & \textit{TBD} \\
2 & \textit{TBD} & \textit{TBD} & \textit{TBD} & \textit{TBD} & \textit{TBD} \\
4 & \textit{TBD} & \textit{TBD} & \textit{TBD} & \textit{TBD} & \textit{TBD} \\
8 & \textit{TBD} & \textit{TBD} & \textit{TBD} & \textit{TBD} & \textit{TBD} \\
\bottomrule
\end{tabular}
\end{table}

% ─────────────────────────────────────────────────────────────────────────────
\section{Analysis}

% TODO: write after results are in. Address:
% - Where speedup is linear (pure VADD on large arrays)
% - Where it deviates (loop overhead: BNE, SADDI are always 1 element per instruction)
% - Fmax vs LANES trend
% - Cost per lane in LUTs

% ─────────────────────────────────────────────────────────────────────────────
\section{Limitations and Future Work}

The current design is non-pipelined (one instruction every 3~cycles). A
3-stage pipeline (Fetch, Decode/Execute, Writeback) would improve throughput
to one instruction per cycle with no data hazards (no register forwarding needed
for the vector path, since the FSM is single-issue).

The VMUL stretch goal (low 32 bits of product) is not implemented in Phase~I;
adding it would utilise FPGA DSP48 blocks and allow benchmark~B4.

Cross-lane operations (reductions, shuffles) are not supported in the base design.
Phase~II scope includes evaluating the cost of a tree-reduce over all lanes.

% ─────────────────────────────────────────────────────────────────────────────
\section{Conclusion}

We have designed and verified a parameterized SIMD VPU in synthesizable Verilog,
demonstrating data-level parallelism at LANES\,$\in\{1,2,4,8\}$. The design
faithfully implements Flynn's SIMD model: a single instruction drives all lanes in
the same clock cycle. Verification against a Python golden model gives confidence
in the RTL correctness. Synthesis results quantify the LUT and frequency cost of
wider datapaths, providing an empirical data point for the classic area-performance
trade-off in SIMD hardware design.

% ─────────────────────────────────────────────────────────────────────────────
\bibliographystyle{IEEEtran}
\bibliography{refs}

\end{document}
